%--------------------------------------------------------------
% thesis.tex 
%--------------------------------------------------------------
% Corso di Laurea in Informatica 
% http://if.dsi.unifi.it/
% @Facolt\`a di Scienze Matematiche, Fisiche e Naturali
% @Universit\`a degli Studi di Firenze
%--------------------------------------------------------------
% - template for the main file of Informatica@Unifi Thesis 
% - based on Classic Thesis Style Copyright (C) 2008 
%   Andr\'e Miede http://www.miede.de   
%--------------------------------------------------------------
\documentclass[twoside,openright,titlepage,fleqn,
	headinclude,12pt,a4paper,BCOR5mm,footinclude]{scrbook}
%--------------------------------------------------------------
\newcommand{\myItalianTitle}{LIBRERIA PYTHON PER L'OCR DI DOCUMENTI D'IDENTIT\`A\xspace}
\newcommand{\myEnglishTitle}{PYTHON LIBRARY TO PERFORM OCR ON IDENTITY DOCUMENTS\xspace}
\newcommand{\myDegree}{Corso di Laurea in Informatica\xspace}
\newcommand{\myName}{Alessio Falai\xspace}
\newcommand{\myProf}{Maria Cecilia Verri\xspace}
\newcommand{\myOtherProf}{Correlatore\xspace}
\newcommand{\mySupervisor}{Michele Barbagli\xspace}
\newcommand{\myFaculty}{
	Scuola di Scienze Matematiche, Fisiche e Naturali\xspace}
\newcommand{\myUni}{\protect{
	Universit\`a degli Studi di Firenze}\xspace}
\newcommand{\myLocation}{Firenze\xspace}
\newcommand{\myTime}{Anno Accademico 2018-2019\xspace}
\newcommand{\myVersion}{Version 0.1\xspace}
%--------------------------------------------------------------
\usepackage[italian]{babel}
\usepackage[latin1]{inputenc} 
\usepackage[T1]{fontenc} 
\usepackage[square,numbers]{natbib} 
\usepackage[fleqn]{amsmath}  
\usepackage{ellipsis}
\usepackage{listings}
\usepackage{subfig}
\usepackage{caption}
\usepackage{appendix}
\usepackage{siunitx}
%--------------------------------------------------------------
\usepackage{dia-classicthesis-ldpkg}
%--------------------------------------------------------------
% Options for classicthesis.sty:
% tocaligned eulerchapternumbers drafting linedheaders 
% listsseparated subfig nochapters beramono eulermath parts 
% minionpro pdfspacing
\usepackage[eulerchapternumbers,linedheaders,subfig,beramono,eulermath,
parts]{classicthesis}
%--------------------------------------------------------------
\newlength{\abcd} % for ab..z string length calculation
% how all the floats will be aligned
\newcommand{\myfloatalign}{\centering} 
\setlength{\extrarowheight}{3pt} % increase table row height
\captionsetup{format=hang,font=small}
%--------------------------------------------------------------
% Layout setting
%--------------------------------------------------------------
\usepackage{geometry}
\geometry{
	a4paper,
	ignoremp,
	bindingoffset = 1cm, 
	textwidth     = 13.5cm,
	textheight    = 21.5cm,
	lmargin       = 3.5cm, % left margin
	tmargin       = 4cm    % top margin 
}

\lstset{
  	frame=tb,
	language=Matlab,
  	aboveskip=3mm,
  	belowskip=3mm,
  	showstringspaces=false,
  	columns=flexible,
  	basicstyle={\small\ttfamily},
  	numbers=none,
  	breaklines=true,
  	breakatwhitespace=true,
  	tabsize=3
}
%--------------------------------------------------------------
\begin{document}
\frenchspacing
\raggedbottom
\pagenumbering{roman}
\pagestyle{plain}
%--------------------------------------------------------------
% Frontmatter
%--------------------------------------------------------------
%--------------------------------------------------------------
% title.tex (use thesis.tex as main file)
%--------------------------------------------------------------
\begin{titlepage}
	\begin{center}
   	\large
      \hfill
      \vfill
      \begingroup
         \includegraphics[scale=0.15]{logo/logo-unifi}\\
%			\spacedallcaps{\myUni} \\ 
			\myFaculty \\
			\myDegree \\ 
			\vspace{0.5cm}
         \vspace{0.5cm}    
         Tesi di Laurea    
      \endgroup 
      \vfill 
      \begingroup
      	\color{Maroon}\spacedallcaps{\myItalianTitle} \\ $\ $\\
      	\spacedallcaps{\myEnglishTitle} \\ 	
	\bigskip
      \endgroup
      \spacedlowsmallcaps{\myName}
      \vfill 
      \vfill
      Relatore: \emph\myProf\\
      \vfill
      \vfill
      \myTime
      \vfill                      
	\end{center}        
\end{titlepage}   
%--------------------------------------------------------------
% back titlepage
%--------------------------------------------------------------
   \newpage
	\thispagestyle{empty}
	\hfill
	\vfill
	\noindent\myName: 
	\textit{\myItalianTitle,} 
	\myDegree, \textcopyright\ \myTime
%--------------------------------------------------------------
% back titlepage end
%--------------------------------------------------------------

\pagestyle{scrheadings}
%--------------------------------------------------------------
% Mainmatter
%--------------------------------------------------------------
\pagenumbering{arabic}
% use \cleardoublepage here to avoid problems with pdfbookmark
%\chapter*{Introduzione}
\addcontentsline{toc}{chapter}{Introduzione}
\markboth{Introduzione}{Introduzione}

Questa tesi vuole essere un riassunto delle attivit\`a svolte in un tirocinio curriculare presso l'azienda QI-LAB di Firenze nel periodo \textit{marzo - giugno 2019}. L'offerta di tirocinio prevedeva l'apporto di migliorie, in termini di accuratezza ed efficienza, riguardanti una libreria per il riconoscimento ottico dei caratteri di alcuni campi di documenti d'identit\`a, attualmente solo italiani. La libreria, denominata QI-OCR, era inizialmente il risultato del lavoro compiuto da un precedente tirocinante, \textit{Emilio Cecchini}, che si era occupato dell'implementazione del framework di base, sul quale ho avuto il piacere di lavorare io stesso. QI-OCR viene distribuita come pacchetto \textit{pip} ed \`e stata scritta con linguaggio Python, in versione \textit{3.6.8}, anche se compatibile con la versione minor successiva, ovvero la \textit{3.7.x}, nonch\`e l'ultima versione stabile rilasciata, nella data in cui mi trovo a scrivere questo testo. QI-OCR prevedeva la suddivisione in 4 moduli principali:
\begin{enumerate}
	\item \textbf{Preprocessing}: Si occupa della localizzazione del documento all'interno dell'immagine, del ritaglio dei contorni e del raddrizzamento, tramite l'algoritmo di feature matching \textit{SIFT}.
	\item \textbf{Docfields}: Si occupa di effettuare un ritaglio statico, mediante coordinate predeterminate, dei campi d'interesse del documento gi\`a centrato.
	\item \textbf{OCR}: Si occupa di effettuare una binarizzazione dei campi dell'immagine mediante un algoritmo di sogliatura globale \textit{ad-hoc} e di fornire le immagini prodotte in input a un motore di OCR.
	\item \textbf{Postprocessing}: Si occupa di applicare tecniche di analisi sintattica e semantica sui risultati ritornati dal software OCR.
\end{enumerate}
Il lavoro da me svolto riguarda il miglioramento della libreria QI-OCR, mediante tecniche che verranno descritte successivamente. A tal proposito, la tesi \`e suddivisa in cinque capitoli. Il capitolo \ref{chap:math-prerequisites} presenta una descrizione dei prerequisiti matematici necessari a comprendere correttamente le spiegazioni dei metodi presenti nei capitoli successivi. Il capitolo \ref{chap:image-binarization} riguarda lo studio e l'implementazione di tecniche di binarizzazione e segmentazione di immagini, ovvero soluzioni necessarie per la "pulizia" di una generica immagine contenente testo, che permettano di ottenere idealmente il solo testo di colore nero e tutto il resto di colore bianco. Il capitolo \ref{chap:image-matching} riguarda invece lo studio di algoritmi di segmentazione basata dal riscontro sul modello, in cui vengono introdotti metodi alternativi al ben noto SIFT per l'individuazione della posizione di un documento all'interno di un'immagine. Il capitolo \ref{chap:text-detection} introduce l'applicazione di tecniche avanzate di \textit{text-detection} con lo scopo di colmare problemi derivanti da documenti che non rispettano alcun tipo di template prefissato, come ad esempio la vecchia carta d'identit\`a cartacea italiana, che risulta comunque rilasciabile, in alternativa alla nuova carta d'identit\`a elettronica italiana, in casi di estrema esigenza o per tutti i cittadini che non possono recarsi per motivi di handicap presso il municipio. Infine, il capitolo \ref{chap:improvements} descrive ulteriori migliorie e innovazioni introdotte, sia per quanto riguarda il lato sviluppo/distribuzione (containerizzazione, benchmarks, \dots), sia per quanto riguarda l'esperienza utente (OCR parametrizzabile sui campi del documento, implementazione di diversi motori di OCR, \dots).
 % use \myChapter command instead of \chapter
\tableofcontents
\listoffigures
\cleardoublepage
\thispagestyle{empty}
\begin{flushright}
\null\vspace{\stretch {1}}
\emph{"If a machine is expected to be infallible, it cannot also be intelligent" \break --- Alan Turing} \vspace{\stretch{2}}\null
\end{flushright}
\cleardoublepage
\addcontentsline{toc}{chapter}{Bibliografia}
\markboth{Bibliografia}{Bibliografia}
\begin{thebibliography}{99}

\bibitem{bib:digital-image-processing}{Wilhelm Burger, Mark J. Burge - \emph{Digital Image Processing - An Algorithmic Introduction Using Java, Second Edition}}

\bibitem{bib:convolution}{Raffaele Cappelli - \emph{Fondamenti di Elaborazione di Immagini - Operazioni sulle immagini} - Universit\`a degli Studi di Bologna}

\bibitem{bib:wikipedia}{Wikipedia - \emph{\url{https://it.wikipedia.org}}}

\bibitem{bib:top-hat-paper}{Guocheng Wang, Yiwen Wang, Hui Li, Xuanqi Chen, Haitao Lu,
Yanpeng Ma, Chun Peng, Yijun Wang, Linyao Tang - \emph{Morphological Background Detection and Illumination Normalization of Text Image with Poor Lighting} - PLOS ONE}

\bibitem{bib:opencv}{OpenCV documentation - \emph{\url{https://docs.opencv.org}}}

\bibitem{bib:skin-thesis}{Simone Parca - \emph{Studio e sviluppo di algoritmi per l'analisi di immagini della pelle acquisite tramite un sensore capacitivo} - Universit\`a degli Studi di Bologna}

\bibitem{bib:binary-images-connectivity}{Francesco Tortorella - \emph{Elementi di geometria delle immagini digitali binarie} - Universit\`a degli studi di Cassino e del Lazio Meridionale}

\bibitem{bib:advanced-morphological}{\emph{Advanced morphological image processing} - University of Groningen}

\bibitem{bib:niblack}{W. Niblack - \emph{An introduction to Digital Image Processing} - Prentice-Hall}

\bibitem{bib:sauvola}{J. Sauvola, M. Pietikainen - \emph{Adaptive document image binarization} - Pattern Recognition 33(2), pp. 225-236}

\bibitem{bib:medium-feature-matching}{Deepanshu Tyagi - \emph{Introduction To Feature Detection And Matching} - \url{https://medium.com}}

\bibitem{bib:text-extraction}{P. Nagabhushan, S. Nirmala - \emph{Text Extraction in Complex Color Document Images for Enhanced Readability}}

\bibitem{bib:sift}{David G. Lowe - \emph{Distinctive Image Features from Scale-Invariant Keypoints} - University of British Columbia}

\bibitem{bib:surf}{Herbert Bay, Tinne Tuytelaars, Luc Van Gool - \emph{SURF: Speeded Up Robust Features}}

\bibitem{bib:kaze}{Pablo Fern\`andez Alcantarilla, Adrien Bartoli, Andrew J. Davison - \emph{KAZE Features}}

\bibitem{bib:akaze}{Pablo F. Alcantarilla, Jes\`us Nuevo, Adrien Bartoli - \emph{Fast Explicit Diffusion for Accelerated Features in Nonlinear Scale Spaces}}

\bibitem{bib:brisk}{Stefan Leutenegger, Margarita Chli, Roland Y. Siegwart - \emph{BRISK: Binary Robust Invariant Scalable Keypoints}}

\bibitem{bib:orb}{Ethan Rublee, Vincent Rabaud, Kurt Konolige, Gary Bradski - \emph{ORB: an efficient alternative to SIFT or SURF}}

\bibitem{bib:prewitt}{J. Prewitt - \emph{Object enhancement and extraction} - Picture Processing and Psychopictorics, pp. 415-431}

\bibitem{bib:canny}{John Canny - \emph{A Computational Approach to Edge Detection}}

\bibitem{bib:hed}{Saining Xie, Zhuowen Tu - \emph{Holistically-Nested Edge Detection}}

\bibitem{bib:mser}{J. Matas, O. Chum, M. Urban, T. Pajdla - \emph{Robust Wide Baseline Stereo from Maximally Stable Extremal Regions}}

\bibitem{bib:mser-linear}{David Nist\`er, Henrik Stew\`enius - \emph{Linear Time Maximally Stable Extremal Regions}}

\bibitem{bib:connected-components-fast}{Azriel Rosenfeld, John L. Pfaltz - \emph{Sequential Operations in Digital Picture Processing}}

\bibitem{bib:mser-canny}{Huizhong Chen, Sam S. Tsai, Georg Schroth, David M. Chen, Radek Grzeszczuk and Bernd Girod - \emph{Robust text detection in natural images with edge-enhanced Maximally Stable Extremal Regions}}

\bibitem{bib:east}{Xinyu Zhou, Cong Yao, He Wen, Yuzhi Wang, Shuchang Zhou, Weiran He, Jiajun Liang - \emph{EAST: An Efficient and Accurate Scene Text Detector}}

\end{thebibliography}

%--------------------------------------------------------------
\end{document}
%--------------------------------------------------------------
